\chapter{Proposed architecture}
\label{chap:proposed-architecture}

This chapter presents the three components of the learning system: the agent of each \ac{V2V} pair, the
same in every method (Section~\ref{sec:pa-agent}); how the agents share what they learn, through \ac{FedAvg}
at the \ac{BS} (Section~\ref{sec:pa-fedavg}) or through the gossip aggregation that this thesis proposes
(Section~\ref{sec:pa-gossip}); and the reward from which each agent learns, shared or individual
(Section~\ref{sec:pa-individual-reward}). It closes with the resulting methods
(Section~\ref{sec:pa-methods}) and the communication cost of each form of aggregation
(Section~\ref{sec:pa-comm-cost}).

%%
\section{The agent of each pair}
% El agente D3QN de cada par: la red (dos capas ocultas de 128 neuronas y cabezas de duelo), la red objetivo, el buffer con PER, la exploración y el entrenamiento.
\label{sec:pa-agent}

Each \ac{V2V} pair has its own \ac{D3QN} agent (Section~\ref{sec:bg-rl}) with two neural networks of the
same shape: the \termdef{online network}, with weights $\vtheta_k$, which decides and learns, and the target
network, with weights $\vtheta_k^{-}$, which is used to compute the learning targets.

\paragraph{The network.} The network takes the observation of 46 numbers (Equation~\ref{eq:sm-observation}),
passes it through two hidden layers of 128 neurons with ReLU activation, and splits into two outputs
(Figure~\ref{fig:pa-network}): the value of the state, $V(s)$, and the advantage of each of the 60 actions,
$A(s, a)$. The dueling architecture \cite{wang+16dueling} combines them as
\begin{equation}
  Q(s, a) = V(s) + A(s, a) - \frac{1}{60} \sum_{a'} A(s, a')
  \label{eq:pa-dueling}
\end{equation}
where subtracting the mean advantage makes the split between value and advantages unique. (The advantage
$A(s, a)$ should not be confused with the penalty $A$ of Equation~\ref{eq:sm-utility}.) The network has
\num{30397} weights, which take \num{121588} bytes in single precision, about \SI{0.12}{\mega\byte}. That is
the size of a model each time it is transferred.

\begin{figure}[tbp]
  \centering
  \begin{tikzpicture}[execute at begin picture={\linespread{1}\selectfont},
      layer/.style={draw, thick, rounded corners=2pt, minimum height=8.5mm, minimum width=11.5mm,
        align=center, font=\scriptsize, inner sep=1.5pt},
      arrow/.style={-{Stealth[length=1.6mm]}, thick}]
    \node[layer, fill=gray!10] (in) {Input\\46};
    \node[layer, fill=cfedavg!12, right=2.8mm of in] (h1) {Hidden\\128};
    \node[layer, fill=cfedavg!12, right=2.8mm of h1] (h2) {Hidden\\128};
    \node[layer, fill=cgossip!18, above right=0.6mm and 4mm of h2] (v) {Value\\1};
    \node[layer, fill=cgossip!18, below right=0.6mm and 4mm of h2] (a) {Advantages\\60};
    \node[layer, fill=cumbral!12, right=17.5mm of h2] (q) {$Q$\\60};
    \draw[arrow] (in) -- (h1);
    \draw[arrow] (h1) -- (h2);
    \draw[arrow] (h2) -- (v);
    \draw[arrow] (h2) -- (a);
    \draw[arrow] (v) -- (q);
    \draw[arrow] (a) -- (q);
  \end{tikzpicture}
  \caption[The network of each agent]{The neural network of each agent, with the number of values in each
    layer. The hidden layers use ReLU activation. The two outputs are combined according to
    Equation~\ref{eq:pa-dueling}, and the chosen action is the one with the largest $Q$.}
  \label{fig:pa-network}
\end{figure}

\paragraph{The memory.} Each agent stores its transitions
$(s_k^t,\allowbreak a_k^t,\allowbreak r_k^t,\allowbreak s_k^{t+1})$ in its own
prioritized replay buffer \cite{schaul+16per}, which holds \num{10000} transitions, that is, the last 50
episodes. Transition $i$ is drawn with probability
\begin{equation}
  P(i) = \frac{p_i^{\omega}}{\sum_j p_j^{\omega}}, \qquad p_i = |\delta_i| + \epsilon_{\mathrm{p}}
  \label{eq:pa-per}
\end{equation}
\begin{wherelist}
  \item $\delta_i$ is the last \ac{TD} error of the transition, and $\epsilon_{\mathrm{p}} = 10^{-5}$ keeps
    every transition drawable;
  \item $\omega = 0.6$ sets how strongly the priorities count; with $\omega = 0$, sampling is uniform.
\end{wherelist}
A new transition enters with the highest priority in the buffer, so that it is used at least once. To
correct the bias of this sampling, each drawn transition is weighted by
$w_i \propto (N_{\mathcal{B}}\, P(i))^{-\beta}$, normalized so that the largest weight in the buffer is 1,
where $N_{\mathcal{B}}$ is the number of stored transitions. The exponent $\beta$ grows linearly from 0.4 to 1
over the training.

\paragraph{Learning.} In each slot, after storing its transition, each agent draws a minibatch $\mathcal{D}$
of 128 transitions from its buffer and takes one step of the Adam optimizer, with a learning rate of
$10^{-3}$, on the loss
\begin{equation}
  \mathcal{L}(\vtheta_k) = \frac{1}{|\mathcal{D}|} \sum_{i \in \mathcal{D}} w_i\,
    \ell\bigl( y_i - Q(s_i, a_i; \vtheta_k) \bigr)
  \label{eq:pa-loss}
\end{equation}
\begin{wherelist}
  \item $y_i$ is the target of double \ac{DQN} \cite{vanhasselt+16double}: the online network chooses the
    next action and the target network evaluates it,
    \begin{equation}
      y_i = r_i + \gamma\, Q\bigl(s_i',\, \argmax_{a'} Q(s_i', a'; \vtheta_k);\; \vtheta_k^{-}\bigr)
      \label{eq:pa-target}
    \end{equation}
    with discount factor $\gamma = 0.9$;
  \item $\ell$ is the Huber loss \cite{huber64robust}: $\ell(x) = x^2/2$ if $|x| \le 1$, and
    $\ell(x) = |x| - 1/2$ otherwise. It grows like the squared error for small errors and only linearly for
    large ones.
\end{wherelist}
Learning starts as soon as the buffer holds 128 transitions. Every 400 learning steps, the target network
copies the weights of the online network. The episodes are cut at 200 slots, but the problem has no natural
end, so the target always includes the value of the next state.

\paragraph{Exploration.} The agent follows an $\varepsilon$-greedy policy (Section~\ref{sec:bg-rl}). The
value of $\varepsilon$ falls linearly from 1 to 0.05 during the first 60\% of training, that is, until
episode 1200 of 2000, and then stays at 0.05.

Section~\ref{sec:exp-differences} lists which of these choices differ from those of \citet{li+22fedmarl},
and Appendix~\ref{apx:diagnosis} explains how the discount factor was chosen.
Algorithm~\ref{alg:pa-training} sums up the whole training, which is the same for every method except in
the aggregation step.

\begin{algorithm}[tbp]
  \caption{Training of the $K$ agents with periodic aggregation}
  \label{alg:pa-training}
  \begin{algorithmic}[1]
    \Require form of aggregation (\acs{FedAvg}, gossip, or none); reward from which the agents learn
      (shared or individual)
    \State For each agent $k$: initialize $\vtheta_k$ at random, set $\vtheta_k^{-} \gets \vtheta_k$, and
      create an empty buffer $\mathcal{B}_k$
    \For{each episode $e = 1, \dots, 2000$}
      \State Draw the positions and find the neighbors $\Nset_k$ of each pair
      \For{each slot $t = 1, \dots, 200$}
        \State Each agent $k$ observes $s_k^t$ and picks $a_k^t$ with its $\varepsilon$-greedy policy
        \State The environment applies the actions and gives each agent $r_k^t$ and $s_k^{t+1}$
        \For{each agent $k$}
          \State Store $(s_k^t, a_k^t, r_k^t, s_k^{t+1})$ in $\mathcal{B}_k$ with the highest priority
          \State Draw 128 transitions from $\mathcal{B}_k$ and take one Adam step on
            Equation~\ref{eq:pa-loss}
          \State Every 400 learning steps: $\vtheta_k^{-} \gets \vtheta_k$
        \EndFor
        \If{50 slots have passed since the last round}
          \State Aggregate the online networks: \acs{FedAvg}, gossip (Algorithm~\ref{alg:pa-gossip}), or
            nothing
        \EndIf
        \State Lower $\varepsilon$ according to the exploration schedule
      \EndFor
    \EndFor
  \end{algorithmic}
\end{algorithm}

%%
\section{Centralized aggregation: FedAvg}
% El esquema de referencia: cada 50 slots la BS recibe los pesos de todos los agentes, los promedia ponderados por su experiencia y les devuelve el promedio.
\label{sec:pa-fedavg}

In the scheme of \citet{li+22fedmarl}, which this thesis uses as the reference, the \ac{BS} acts as the
federated server. Every $N_{\mathrm{a}} = 50$ slots, that is, four times per episode, there is a round:
\begin{enumerate}[nosep]
  \item each agent sends the weights of its online network to the \ac{BS};
  \item the \ac{BS} averages them with Equation~\ref{eq:bg-fedavg}, where $n_k$ is the number of transitions
    with which agent $k$ trained since the last round;
  \item the \ac{BS} sends the average $\vtheta_{\mathrm{g}}$ back to all the agents, and each one replaces
    its online network with it.
\end{enumerate}
The target networks are not shared: each agent keeps its own. Since every agent trains on 128 transitions in
every slot, all the $n_k$ are equal, and the average is a simple one. After each round, all the agents have
exactly the same online network; between one round and the next, each one adjusts it with its own
experience.

%%
\section{Decentralized aggregation by gossip}
% La propuesta: en las mismas rondas, cada agente promedia su red solo con las de sus vecinos a 300 m o menos, sin pasar por la BS.
\label{sec:pa-gossip}

The aggregation that this thesis proposes takes the \ac{BS} out of learning. The rounds take place at the
same instants as in \ac{FedAvg}, but each agent $k$ averages its network only with those of its neighbors
$\Nset_k$, the pairs whose transmitters are at most \SI{300}{\meter} away
(Equation~\ref{eq:sm-neighbors}). In each round,
\begin{equation}
  \vtheta_k \leftarrow \sum_{j \in \mathcal{G}_k} \frac{n_j}{\sum_{i \in \mathcal{G}_k} n_i}\;
    \bar{\vtheta}_j , \qquad \mathcal{G}_k = \{k\} \cup \Nset_k
  \label{eq:pa-gossip}
\end{equation}
\begin{wherelist}
  \item $\bar{\vtheta}_j$ are the weights of the online network of agent $j$ just before the round;
  \item $\mathcal{G}_k$ is the group formed by agent $k$ and its neighbors;
  \item $n_j$ is, as in \ac{FedAvg}, the number of transitions with which agent $j$ trained since the last
    round. Since all the $n_j$ are equal, each agent ends up with the simple average of its network and those
    of its neighbors, the gossip rule of Equation~\ref{eq:bg-gossip}.
\end{wherelist}
Three details complete the definition. The rounds are synchronous: all the agents average the weights from
before the round, so the result does not depend on the order in which they are processed. An agent with no
neighbors keeps its network and transfers nothing. And, as in \ac{FedAvg}, the target networks are not
shared. Algorithm~\ref{alg:pa-gossip} sums up one round, and Figure~\ref{fig:bg-fedavg-gossip} compares the
two forms of aggregation.

\begin{algorithm}[tbp]
  \caption{One round of gossip aggregation}
  \label{alg:pa-gossip}
  \begin{algorithmic}[1]
    \Require online networks $\vtheta_1, \dots, \vtheta_K$; transitions $n_k$ trained on since the last
      round; neighbors $\Nset_k$; optional cap $S$
    \State $\bar{\vtheta}_k \gets \vtheta_k$ for every $k$ \Comment{all average the old weights}
    \For{each agent $k$}
      \State $\mathcal{S}_k \gets \Nset_k$
      \If{$S$ is set and $|\Nset_k| > S$}
        \State $\mathcal{S}_k \gets$ $S$ neighbors drawn from $\Nset_k$ at random, without replacement
      \EndIf
      \If{$\mathcal{S}_k \ne \emptyset$}
        \State Receive $\bar{\vtheta}_j$ from each $j \in \mathcal{S}_k$ \Comment{one transfer each}
        \State $\mathcal{G} \gets \{k\} \cup \mathcal{S}_k$
        \State $\textstyle\vtheta_k \gets \sum_{j \in \mathcal{G}} \bigl(n_j / \textstyle\sum_{i \in \mathcal{G}}
          n_i\bigr)\, \bar{\vtheta}_j$
      \EndIf
    \EndFor
    \State $n_k \gets 0$ for every $k$
  \end{algorithmic}
\end{algorithm}

\subsection{Why averaging with neighbors makes sense}
% Todos los pares tienen la misma observación, las mismas acciones y la misma forma de recompensa, así que una misma red les sirve a todos y compartir pesos transfiere experiencia entre vecinos.
\label{sec:pa-gossip-why}

Averaging the networks of different agents only makes sense if one network can serve all of them. Here it
can: all the pairs have the same kind of observation, the same actions, and the same form of reward, and their
positions are drawn in the same way. A policy that decides from the local observation serves, in principle,
any pair. Sharing weights therefore gives each agent what the others learned in situations it did not live
through, as in horizontal federated reinforcement learning (Section~\ref{sec:bg-fl}).

The schemes differ in how close they bring the networks to each other. If $\bm{\Theta}$ is the matrix whose
rows are the weights of the $K$ agents, a round of any of the three schemes can be written as
$\bm{\Theta} \leftarrow \bm{W} \bm{\Theta}$, where row $k$ of $\bm{W}$ says with what weight each agent enters
the new model of $k$:
\begin{itemize}[nosep]
  \item without aggregation (\ac{DTE}), $\bm{W}$ is the identity: nobody changes;
  \item in \ac{FedAvg}, every entry of $\bm{W}$ is $1/K$: all the agents become equal;
  \item in gossip, row $k$ is $1/(1 + |\Nset_k|)$ in the columns of $k$ and its neighbors, and zero
    elsewhere: each agent comes closer to its neighbors, but not to the others.
\end{itemize}
Gossip thus lies between learning alone and forcing everyone to be equal. In one round, knowledge advances
only one neighbor; but since the graph changes in every episode, over many rounds what each agent knows ends
up reaching all the others (Section~\ref{sec:bg-dfl}). We call this aggregation \termdef{dynamic} because the
groups that average are not fixed: the positions of the vehicles decide them anew in every episode.

The aggregation needs no new information either. The neighborhood it uses is the same \SI{300}{\meter} one
that already appears in the observation of each agent (Equation~\ref{eq:sm-observation}), and the models
travel over the same direct link between vehicles that, according to \citet{li+22fedmarl}, already carries
the channel choices of the neighbors.

\subsection{Variant with at most \texorpdfstring{$S$}{S} neighbors}
% Promediar en cada ronda con a lo sumo S vecinos elegidos al azar (siguiendo a Ge et al) para acotar la comunicación.
\label{sec:pa-gossip-s}

With many pairs, each vehicle has many neighbors, and communication grows fast (Section~\ref{sec:pa-comm-cost}).
Following \citet{ge+24dfl}, the capped variant sets a maximum number $S$ of neighbors per round: if an agent
has more than $S$ neighbors, it draws $S$ of them at random, without replacement, and averages only with
those. The draw is repeated in every round, so over the rounds the agent ends up averaging with all its
neighbors. With the cap, the relation is no longer symmetric: that agent $k$ takes the model of $j$ does not
mean that $j$ takes the model of $k$. The experiments use $S = 2$ and $S = 3$.

%%
\section{Individual reward}
% La recompensa propia de cada agente (su cola, la confiabilidad de su enlace y el CUE del canal que eligió) en lugar de la suma de todos, y qué alineación se sacrifica a cambio de que sea más fácil de aprender.
\label{sec:pa-individual-reward}

The shared reward (Equation~\ref{eq:sm-shared-reward}) adds up $M + 2K$ terms: one for each \ac{CUE} and two
for each pair. The action of one agent directly changes very few of them: the term of the \ac{CUE} of the
channel it chose, those of its own queue and its own reliability, and, if it shares the channel with other
pairs, the reliability of those pairs. The rest of the sum depends on what the others do. From the point of
view of each agent, that part is noise, and it grows with $K$: with 30 pairs, the shared reward has 75 terms.
This is the problem of learnability described in Section~\ref{sec:bg-marl}.

The individual reward that this thesis proposes judges each agent only by the terms that its action touches.
The reward of agent $k$ in slot $t$ is
\begin{equation}
  \begin{aligned}
    r_k^t = {} & \lambda_1\, U\Bigl( \frac{R_{m_k^t}^{\mathrm{c},t} - \Rmin}{\Rmin} \Bigr)
      + \lambda_2\, U\Bigl( \frac{\Qmax - Q_k^{t+1}}{\Qmax} \Bigr) \\
      & + \lambda_3\, U\Bigl( \frac{\mu_k^t}{\gth} \Bigr)
  \end{aligned}
  \label{eq:pa-individual-reward}
\end{equation}
\begin{wherelist}
  \item the first term looks only at the \ac{CUE} of the channel $m_k^t$ that the agent chose, the user that
    its transmission can harm. All the pairs on the same channel share this term;
  \item the second looks only at its own queue, and measures it at the \emph{end} of the slot, $Q_k^{t+1}$
    (Equation~\ref{eq:sm-queue}), instead of at its start;
  \item the third looks only at its own reliability margin (Equation~\ref{eq:sm-margin});
  \item $U$, $\lambda_1$, $\lambda_2$, $\lambda_3$, and $A$ are those of the shared reward.
\end{wherelist}

\paragraph{Why the queue at the end of the slot.} In the shared reward, the queue term uses $Q_k^t$, the queue
with which the slot \emph{starts}, which the action of the slot can no longer change. If a pair stays silent
in the slot in which its queue crosses the maximum, the penalty arrives one slot later, mixed with the
consequences of another action. With $Q_k^{t+1}$, the penalty falls in the same slot as the action that
caused it.

\paragraph{The example, with the individual reward.} In the tiny network of
Section~\ref{sec:sm-shared-reward}, both pairs received $-5.2$. With the individual reward, pair A receives
$1.2 + 0.4 + 0.2 = 1.8$, and pair B receives $1.2 - 5 - 2 = -5.8$: each one receives the consequence of what it
did. The queue at the end of the slot adds another difference. If pair B started the slot with \num{99500}
bits, its queue term in the shared reward is $0.5 \times U(0.005) = 0.0025$: the pair stays silent, the queue
overflows, and, in that slot, there is no penalty. With the individual reward, the queue at the end of the
slot is \num{100500} bits, and the penalty of $-5$ arrives in the slot of the silence.

\paragraph{What is gained and what is lost.} In the terms of Section~\ref{sec:bg-marl}, the individual reward
is much more learnable than the shared one, because almost everything that moves it depends on the agent
itself. In exchange, it is less aligned with the objective of the system: it does not count the interference
that the agent causes to the other pairs on its channel, which the shared reward does count, through the
reliability of those pairs. An agent could improve its reward by transmitting hard at the expense of the
pairs on its channel. Whether the trade is worth it is an empirical question, which
Section~\ref{sec:res-reward} answers.

\paragraph{What information it needs.} The shared reward requires the \ac{BS} to know, in every slot, the queue
and the margin of every pair. The individual reward uses the agent's own queue, its own margin, which its
receiver helps it estimate, and the rate of the \ac{CUE} of its channel, which the \ac{BS} already knows and
would have to broadcast. It is not completely local, but the vehicles no longer report anything to the
\ac{BS}. It also changes only what the agents learn from: every method is still measured with the shared
reward (Equation~\ref{eq:sm-shared-reward}).

%%
\section{Resulting methods}
% Tabla de combinaciones de agregación y recompensa: el esquema original (FedAvg con recompensa compartida), la propuesta (gossip con recompensa individual) y las demás, que separan el efecto de cada cambio.
\label{sec:pa-methods}

Combining the form of aggregation with the reward gives the learning methods that this thesis compares
(Table~\ref{tab:pa-methods}). The original scheme of \citet{li+22fedmarl} is \ac{FedAvg} with the shared
reward, and the complete proposal is gossip with the individual reward. The other combinations separate the
effect of each change.

\begin{table}[tbp]
  \centering
  \footnotesize
  \caption[Combinations of aggregation and reward]{Combinations of aggregation and reward that are
    evaluated. All of them are measured with the shared reward of Equation~\ref{eq:sm-shared-reward}.}
  \label{tab:pa-methods}
  \begin{tabularx}{\linewidth}{@{}l >{\raggedright\arraybackslash}X >{\raggedright\arraybackslash}X@{}}
    \toprule
    Aggregation & Shared reward & Individual reward \\
    \midrule
    None (\acs{DTE}) & control & control \\
    \acs{FedAvg} at the \acs{BS} & original scheme & new reward only \\
    Gossip, all neighbors & new aggregation only & \textbf{complete proposal} \\
    Gossip, $S = 2$ or 3 & --- & proposal, bounded cost \\
    \bottomrule
  \end{tabularx}
\end{table}

%%
\section{Communication cost}
% Cómo se mide el costo (transferencias de modelos de 121 588 bytes y carga del nodo más cargado) y las fórmulas del costo de FedAvg y de gossip.
\label{sec:pa-comm-cost}

We count the cost of a form of aggregation as the number of \termdef{transfers}: each time the model of one
node reaches another node counts as one. Each transfer moves \num{121588} bytes. Both the total and its
distribution matter: the \termdef{load} of a node is the number of transfers in which it takes part, sending
or receiving, and the busiest node is the first to saturate. Table~\ref{tab:pa-cost} gathers the expressions
derived below.

\paragraph{FedAvg.} In each round, every agent sends its model to the \ac{BS} and receives the average: $2K$
transfers, all of them through the \ac{BS}, whose load grows linearly with the number of pairs.

\paragraph{Gossip with all neighbors.} In each round, every agent receives the model of each of its
neighbors, so the total is the sum of the degrees, $\textstyle\sum_k |\Nset_k|$. Since the neighbor relation is
symmetric, each vehicle sends as many models as it receives, and its load is $2|\Nset_k|$; the busiest one is
the one with the most neighbors. The total depends on how dense the network is. If $p$ is the probability that
two vehicles drawn at random in the cell are at most \SI{300}{\meter} apart, each pair has, on average,
$p(K-1)$ neighbors, and the total per round is about
\begin{equation}
  \sum_{k \in \Kset} |\Nset_k| \approx p\, K (K - 1)
  \label{eq:pa-cost-gossip}
\end{equation}
which grows with the square of $K$. In a cell without edges, $p$ would be $(300/500)^2 = 0.36$; the edge of
the cell lowers it, because vehicles near the edge have fewer possible neighbors. In the simulations,
$p \approx 0.27$ (Section~\ref{sec:res-consensus}). Hence, with few pairs, gossip transfers less than
\ac{FedAvg}, and with many, more; but no node takes part in every transfer.

\paragraph{Gossip with at most $S$ neighbors.} Each agent receives at most $S$ models per round, so the total
does not exceed $SK$: with $S = 2$, at most $2K$, the same as \ac{FedAvg}. Each vehicle also sends its model
to the neighbors that drew it. On average, a vehicle sends as many models as it receives, so the mean load is
at most $2S$; the busiest vehicle may be drawn by more neighbors, but its load no longer grows with the number
of pairs (Section~\ref{sec:res-communication}).

\begin{table}[tbp]
  \centering
  \small
  \caption[Communication cost of each form of aggregation]{Communication cost per round of each form of
    aggregation. Only \acs{FedAvg} goes through the \acs{BS}. An episode has four rounds, and each transfer
    moves \num{121588} bytes.}
  \label{tab:pa-cost}
  \begin{tabularx}{\linewidth}{@{}l l >{\raggedright\arraybackslash}X@{}}
    \toprule
    Aggregation & Transfers & Busiest node (load) \\
    \midrule
    None (\acs{DTE}) & 0 & --- \\
    \acs{FedAvg} & $2K$ & the \acs{BS} ($2K$) \\
    Gossip & $\textstyle\sum_k |\Nset_k|$ & the vehicle with the most neighbors
      ($\textstyle 2 \max_k |\Nset_k|$) \\
    Gossip, $S$ & at most $SK$ & a vehicle (mean at most $2S$) \\
    \bottomrule
  \end{tabularx}
\end{table}

Two remarks complete the analysis. First, the count is conservative for gossip: over a wireless medium, a vehicle could broadcast its
model once to all its neighbors, while here each reception counts as one transfer. Second, the transfers of
\ac{FedAvg} use the links to the cellular network, and those of gossip the direct link between vehicles.
Section~\ref{sec:res-communication} measures these costs in the simulations.
